
\section{Validation program: does the estimator itself hold up?}
\label{sec:validation}
The headline census numbers mean little unless the estimator fails loudly
when it should and behaves predictably when the data change. Four
experiments test exactly that. All run on a seeded every-second-image
subsample of the training partition ($n=11023$) with reduced folds and
epochs for compute feasibility; the canonical census in
Sec.~\ref{sec:results} uses the full partition, so rates here are
protocol-internal comparisons, not replacements for the headline numbers.
Raw per-fold probability checkpoints and per-experiment JSONs are committed
under \texttt{results/malaria/validation/}.

\subsection{Negative control: the gate refuses random labels}
We shuffle the training labels (seed 42) and rerun the census protocol
verbatim. Out-of-fold accuracy collapses to 49.23\%
--- \emph{below} the 50.51\% majority rate ---
so the admissibility gate \textbf{refuses to emit a noise estimate at
all}. Without the gate, the same machinery would have reported a
``50.72\% label-noise rate'' on pure
noise: confident learning applied to a dataset with no signal fabricates a
census. This experiment is why the gate exists, and it is the control every
label-noise audit of a benchmark should publish alongside its flags.

\subsection{Dose-response: injected flips are recovered monotonically}
We inject synthetic label flips at 2\%, 5\%, 10\% and 20\% on top of the
real labels and rerun the census. The estimated rate rises monotonically
and tracks the true total (real $\approx$11\% $+$ injected), and recall
against the injected set rises with dose (Table~\ref{tab:recovery}). The
estimator is dose-sensitive and directionally calibrated: it is not
saturating on a fixed suspicious subset.
\begin{table}[h]
\centering\small
\caption{Dose-response recovery. Precision/recall are measured against the
injected flips only; overlap with the canonical 249 flags stays
$\approx$90--99 throughout, i.e.\ the real flags persist under
contamination.}
\label{tab:recovery}
\begin{tabular}{r r r r r}
\hline
injected & estimated rate & flags & precision & recall \\ \hline
2\% & 10.96\% & 335 & 0.534 & 0.814 \\
5\% & 12.48\% & 655 & 0.728 & 0.866 \\
10\% & 15.45\% & 1199 & 0.825 & 0.897 \\
20\% & 23.12\% & 2273 & 0.894 & 0.922 \\
\hline
\end{tabular}
\end{table}

\subsection{Sensitivity to protocol choices}
Varying folds (3/5), epochs (4/6) and seed (7/42/123) moves the estimated
rate within a 9.90--11.48\% band around the canonical
11.44\%, and per-run flag sets overlap the canonical 249 by
67--78 images. The rate is stable to protocol choices; the identity of the
marginal flags is not, which is why we report flags as candidates with a
consensus core rather than as verdicts.
\begin{table}[h]
\centering\small
\caption{Protocol sensitivity of the CNN census (subsample protocol).}
\label{tab:sensitivity}
\begin{tabular}{r r r r r r}
\hline
folds & epochs & seed & rate & flags & overlap w.\ 249 \\ \hline
5 & 4 & 42 & 9.90\% & 128 & 75 \\
3 & 4 & 7 & 10.61\% & 115 & 78 \\
3 & 4 & 123 & 10.99\% & 114 & 67 \\
3 & 6 & 42 & 11.48\% & 122 & 78 \\
\hline
\end{tabular}
\end{table}

\subsection{Cross-architecture stability: GCN census}
A GCN-based auditor (same folds, epochs and seed 42 as the CNN sensitivity
runs) estimates 9.59\% with
153 flags. That rate sits 0.31 percentage points
\emph{below} the CNN band floor of 9.90\% --- close to the band, but
not inside it, and we report it that way rather than rounding the
difference away. Flag identity agrees only partially:
85 of the GCN flags coincide with the canonical
CNN 249 (Jaccard 0.268). The honest reading:
the \emph{rate} is architecture-stable, the \emph{marginal flag
identity} is representation-dependent, and the
85-image intersection is a high-confidence
consensus core that both architectures flag independently. A defensible
release is the consensus core plus architecture-labelled candidate tiers,
not a single flat list.

\subsection{Calibration of the deployed probabilities}
\label{sec:calibration}
A diagnostic score is only useful if its probabilities mean what they say.
On the untouched test partition (committed probability dumps, no
retraining), the CNN auditor scores ECE 0.0204 and
Brier 0.0302; the GCN scores ECE 0.0088 and Brier
0.0315 --- both far better calibrated than a base-rate prior
(Brier skill 0.88/0.87). Two independent ECE
implementations and two Brier implementations agree to $10^{-9}$
(\texttt{results/calibration.json}). The clinical reading: a predicted
0.9 is right about nine times in ten on average; calibration is a
population property of the score, not a per-image guarantee, and we report
it as such.

\subsection{Error enrichment: do the flags live where models struggle?}
\label{sec:enrichment}
Per-image out-of-fold loss tells the same story twice. The CNN's own 249
flags are 10.0$\times$ enriched in its top-decile OOF-loss
stratum --- partly by construction, since the flags derive from those same
probabilities, so we report it as a sanity check only. The non-circular
check is cross-architectural: the 68 images flagged \emph{only}
by the GCN auditor (no CNN input) make up
0.31\% of the training set but
3.04\% of the CNN's top-decile loss --- a
9.8$\times$ enrichment, with mean OOF loss
2.02 against 0.09 for unflagged images
(\texttt{results/malaria/error\_enrichment.json}). Images one architecture
suspects genuinely give the other architecture trouble: the flags track a
real, representation-independent difficult stratum, which is what
directional label contamination looks like before expert adjudication.

\subsection{The estimand, stated precisely}
\label{sec:estimand}
The headline 11.44\% is not ``the fraction of wrong labels''. It is the
confident-learning estimator's estimate of the fraction of training
examples whose observed label disagrees with the latent-label posterior
under the specified three-fold out-of-fold probability model:
\begin{equation}
\hat{\pi}_{\mathrm{noise}} \;=\;
\frac{1}{N}\sum_{c\ne c'} \hat{Q}_{c,c'} ,
\end{equation}
where $\hat{Q}$ is the calibrated confident joint over observed class
$c$ and latent class $c'$ estimated from the OOF probabilities
(\texttt{results/malaria/label\_noise\_census.json}). It is an
estimator-derived quantity, not independently verified biological ground
truth; the dose-response experiment bounds its behaviour under known
corruption, and the band bounds its protocol sensitivity.

\subsection{Patient-clustered uncertainty}
\label{sec:clusterboot}
Cells from the same smear are not independent, so image-level intervals
understate uncertainty. Recomputing the headline accuracy intervals with
a percentile cluster bootstrap (resampling unit: smear/patient code, the
filename prefix before \texttt{\_IMG}; 201 clusters, median 9 cells, $B=10{,}000$)
widens the intervals exactly as it should: CNN [95.22, 96.89]\%
vs Wilson [95.29, 96.75]\%, GCN [95.01, 96.81]\%
vs Wilson [95.18, 96.65]\%
(\texttt{results/malaria/cluster\_bootstrap.json}). Both constructions
are reported; the clustered interval is the conservative one.

\subsection{What the method proves, and what it does not}
\label{sec:proves}
\textbf{Supported by the evidence in this paper:} the estimator responds
monotonically to known injected noise; the gate refuses to report under
model collapse; the population estimate is stable across tested protocol
configurations; a similar noise magnitude appears under an independent
auditor architecture; individual candidate identity is
representation-dependent while an 85-image consensus core reproduces;
candidates are cross-architecturally enriched in out-of-fold loss; the
same gated pipeline ports to two further biomedical collections and
refuses an inadmissible case.
\textbf{Not established:} that any individual flagged image is
biologically mislabelled; that 11.44\% is the exact biological
mislabelling fraction; universal generalisation to all biomedical
datasets; or a causal claim about \emph{why} the asymmetry is
directional. Confirming individual candidates requires blinded expert
adjudication, scoped as follow-up.
